% GENERATED FILE. Edit canonical lessons, then run npm run generate:books.
% canonical-source-hash: fnv1a64:9af9a07dbb7430ec
% canonical-lessons: RU-C126-last, RU-C126-famous, RU-C126-important, RU-C126-interesting, RU-C126-boring

\chapter{Last, Famous, Important, Interesting, Boring}
\label{ch:ru-last-famous-important-interesting-boring}
\hlchaptermodality{126}

\begin{chapteropening}
\textbf{By the end of this chapter:} \emph{I can say last, famous, important, interesting and boring.}

Complete the last lesson of chapter 126: I can say last, famous, important, interesting and boring.
\end{chapteropening}

\section[\texorpdfstring{\ru{последний} (poslédniy)}{последний (poslédniy)}]{\ru{последний} (poslédniy) — last}

\label{lesson:RU-C126-last}



\begin{quote}
Before the new one: say the Russian for quiet, then the Russian for noisy.
\end{quote}

\subsection*{You'll want to know: \ru{последний}}
\textbf{\ru{последний}} (\emph{poslédniy}) — \textquotedblleft{}last\textquotedblright{}.

\begin{cousinweb}
\textbf{\ru{последний}} comes from Old Russian \emph{\ru{последь}}, \textquotedblleft{}afterwards\textquotedblright{}, which goes with \emph{\ru{след}}, \textquotedblleft{}a track\textquotedblright{} (Vasmer): the last is the one at the back of the track.
\end{cousinweb}

\begin{grammarlens}[title={What you've built}]
One more word for this chapter.
\end{grammarlens}

\subsection*{Your turn}
\begin{itemize}
\raggedright
  \item \emph{Say it:} \emph{poslédniy}
  \item \emph{Say it:} \emph{poslédniy}, once more
  \item \emph{Say it:} say \emph{poslédniy}
  \item \emph{Recall:} say the Russian for quiet, then the Russian for noisy, then say \emph{poslédniy} again
\end{itemize}

\begin{tcolorbox}[breakable,colback=teal!4,colframe=teal!35!black,title={Before you move on}]
What does \ru{последний} mean? (\textquotedblleft{}Last\textquotedblright{}.) Say it once more.
\end{tcolorbox}
\section[\texorpdfstring{\ru{известный} (izvéstnyy)}{известный (izvéstnyy)}]{\ru{известный} (izvéstnyy) — famous}

\label{lesson:RU-C126-famous}



\begin{quote}
Before the new one: say the Russian for noisy, then the Russian for last.
\end{quote}

\subsection*{You'll want to know: \ru{известный}}
\textbf{\ru{известный}} (\emph{izvéstnyy}) — \textquotedblleft{}famous\textquotedblright{}.

\begin{cousinweb}
\textbf{\ru{известный}} is built on \emph{\ru{известие}}, \textquotedblleft{}news\textquotedblright{}, \emph{\ru{из}-} with \emph{\ru{весть}}, \textquotedblleft{}a message\textquotedblright{}, from \emph{\ru{ведать}}, \textquotedblleft{}to know\textquotedblright{} (Vasmer): \textquotedblleft{}known\textquotedblright{}. Vasmer compares Sanskrit \emph{vitti-}, \textquotedblleft{}knowledge\textquotedblright{}.
\end{cousinweb}

\begin{grammarlens}[title={What you've built}]
One more word for this chapter.
\end{grammarlens}

\subsection*{Your turn}
\begin{itemize}
\raggedright
  \item \emph{Say it:} \emph{izvéstnyy}
  \item \emph{Say it:} \emph{izvéstnyy}, once more
  \item \emph{Say it:} say \emph{izvéstnyy}
  \item \emph{Recall:} say the Russian for noisy, then the Russian for last, then say \emph{izvéstnyy} again
\end{itemize}

\begin{tcolorbox}[breakable,colback=teal!4,colframe=teal!35!black,title={Before you move on}]
What does \ru{известный} mean? (\textquotedblleft{}Famous\textquotedblright{}.) Say it once more.
\end{tcolorbox}
\section[\texorpdfstring{\ru{важный} (vázhnyy)}{важный (vázhnyy)}]{\ru{важный} (vázhnyy) — important}

\label{lesson:RU-C126-important}



\begin{quote}
Before the new one: say the Russian for last, then the Russian for famous.
\end{quote}

\subsection*{You'll want to know: \ru{важный}}
\textbf{\ru{важный}} (\emph{vázhnyy}) — \textquotedblleft{}important\textquotedblright{}.

\begin{cousinweb}
\textbf{\ru{важный}} apparently came from Polish \emph{ważny}, \textquotedblleft{}weighty\textquotedblright{}, from \emph{waga}, \textquotedblleft{}a scale\textquotedblright{}, from German \emph{Waage}, \textquotedblleft{}scales\textquotedblright{} (Vasmer): important is \textquotedblleft{}weighty\textquotedblright{}.
\end{cousinweb}

\begin{grammarlens}[title={What you've built}]
One more word for this chapter.
\end{grammarlens}

\subsection*{Your turn}
\begin{itemize}
\raggedright
  \item \emph{Say it:} \emph{vázhnyy}
  \item \emph{Say it:} \emph{vázhnyy}, once more
  \item \emph{Say it:} say \emph{vázhnyy}
  \item \emph{Recall:} say the Russian for last, then the Russian for famous, then say \emph{vázhnyy} again
\end{itemize}

\begin{tcolorbox}[breakable,colback=teal!4,colframe=teal!35!black,title={Before you move on}]
What does \ru{важный} mean? (\textquotedblleft{}Important\textquotedblright{}.) Say it once more.
\end{tcolorbox}
\section[\texorpdfstring{\ru{интересный} (interésnyy)}{интересный (interésnyy)}]{\ru{интересный} (interésnyy) — interesting}

\label{lesson:RU-C126-interesting}



\begin{quote}
Before the new one: say the Russian for famous, then the Russian for important.
\end{quote}

\subsection*{You'll want to know: \ru{интересный}}
\textbf{\ru{интересный}} (\emph{interésnyy}) — \textquotedblleft{}interesting\textquotedblright{}.

\begin{cousinweb}
\textbf{\ru{интересный}} is built on \emph{\ru{интерес}}, \textquotedblleft{}interest\textquotedblright{}, recorded from 1703, which came perhaps through Polish \emph{interes} or German \emph{Interesse} from Medieval Latin \emph{interesse}, \textquotedblleft{}to matter\textquotedblright{} (Vasmer).
\end{cousinweb}

\begin{grammarlens}[title={What you've built}]
One more word for this chapter.
\end{grammarlens}

\subsection*{Your turn}
\begin{itemize}
\raggedright
  \item \emph{Say it:} \emph{interésnyy}
  \item \emph{Say it:} \emph{interésnyy}, once more
  \item \emph{Say it:} say \emph{interésnyy}
  \item \emph{Recall:} say the Russian for famous, then the Russian for important, then say \emph{interésnyy} again
\end{itemize}

\begin{tcolorbox}[breakable,colback=teal!4,colframe=teal!35!black,title={Before you move on}]
What does \ru{интересный} mean? (\textquotedblleft{}Interesting\textquotedblright{}.) Say it once more.
\end{tcolorbox}
\section[\texorpdfstring{\ru{скучный} (skúchnyy)}{скучный (skúchnyy)}]{\ru{скучный} (skúchnyy) — boring}

\label{lesson:RU-C126-boring}



\begin{quote}
Before the new one: say the Russian for important, then the Russian for interesting.
\end{quote}

\subsection*{You'll want to know: \ru{скучный}}
\textbf{\ru{скучный}} (\emph{skúchnyy}) — \textquotedblleft{}boring\textquotedblright{}.

\begin{cousinweb}
\textbf{\ru{скучный}} is \emph{\ru{скука}}, \textquotedblleft{}boredom\textquotedblright{}, with \emph{-\ru{ный}}. Vasmer will not separate \emph{\ru{скука}} from \emph{\ru{докучать}}, \textquotedblleft{}to pester\textquotedblright{}, but rejects the proposed outside links, so its origin is open.
\end{cousinweb}

\begin{grammarlens}[title={What you've built}]
One more word for this chapter.
\end{grammarlens}

\subsection*{Your turn}
\begin{itemize}
\raggedright
  \item \emph{Say it:} \emph{skúchnyy}
  \item \emph{Say it:} \emph{skúchnyy}, once more
  \item \emph{Say it:} say \emph{skúchnyy}
  \item \emph{Recall:} say the Russian for important, then the Russian for interesting, then say \emph{skúchnyy} again
\end{itemize}

\begin{tcolorbox}[breakable,colback=teal!4,colframe=teal!35!black,title={Before you move on}]
What does \ru{скучный} mean? (\textquotedblleft{}Boring\textquotedblright{}.) Say it once more.
\end{tcolorbox}
